\documentclass[11pt]{amsart}

\usepackage[T1]{fontenc}
\usepackage{lmodern}
\usepackage{microtype}
\usepackage{amsmath,amssymb,amsthm}
\usepackage[hidelinks]{hyperref}

\hypersetup{
  pdftitle={A Counterexample to the Aggregate-Condition Form of Karp--Sitnik Conjecture 1}
}

\newtheorem{theorem}{Theorem}
\newtheorem{proposition}[theorem]{Proposition}
\newtheorem{corollary}[theorem]{Corollary}

\title{A Counterexample to the Aggregate-Condition Form of
Karp--Sitnik Conjecture 1}
\subjclass[2020]{Primary 33C20; Secondary 26D15}
\keywords{generalized hypergeometric function, Karp--Sitnik conjecture,
aggregate condition, majorization, counterexample}
\date{}

\begin{document}

\begin{abstract}
We consider the nonredundant aggregate-condition form of Conjecture~1 of
Karp and Sitnik, in which the componentwise conditions $b_i>a_i>0$ are
replaced by $\sum_i(b_i-a_i)>0$ while positivity is retained. We give the
exact counterexample
\[
 {}_3F_2\!\left(
 \begin{matrix}2,2,2\\1,5\end{matrix};-\frac12
 \right)
 =1032-2544\log\frac32<\frac{25}{49}.
\]
Thus the aggregate condition is insufficient already for $q=2$, $\sigma=2$,
positive integer parameters, and $\prod_i a_i/b_i<1$. A quadratic expansion
also gives an infinite family of counterexamples for every $\sigma>0$.
\end{abstract}

\maketitle

Let $q\ge1$. For $A=(a_1,\ldots,a_q)$ and $B=(b_1,\ldots,b_q)$ in
$(0,\infty)^q$, set
\[
 H_{\sigma,A,B}(x)
 :=
 {}_{q+1}F_q\!\left(
 \begin{matrix}
 \sigma,a_1,\ldots,a_q\\
 b_1,\ldots,b_q
 \end{matrix};-x
 \right),
 \qquad
 \mu(A,B):=\prod_{i=1}^q\frac{a_i}{b_i}.
\]
Karp and Sitnik proved that for $b_i>a_i>0$, $\sigma\ge1$ and $x>-1$,
\begin{equation}\label{eq:bound}
 (1+\mu(A,B)x)^{-\sigma}<H_{\sigma,A,B}(x),
\end{equation}
\cite[Theorem~3]{KS}. Their Conjecture~1 reads, in full, ``Theorem~3 is
true for all $\sigma>0$ and $\sum_{i=1}^{q}(b_i-a_i)>0$''
\cite[Conjecture~1]{KS}. We read the condition
\begin{equation}\label{eq:aggregate}
 \sum_{i=1}^q(b_i-a_i)>0
\end{equation}
as replacing the componentwise inequalities $b_i>a_i$ of Theorem~3, all
parameters remaining positive. Three features of \cite{KS} force this
reading. First, \eqref{eq:aggregate} is automatic under $b_i>a_i$, so on
the componentwise reading it would be an idle clause. Second, Conjecture~2
of \cite{KS} lists ``$0<\sigma\le\min(a_1,a_2,\ldots,a_q)$,
$\sum_{i=1}^{q}(b_i-a_i)>0$ and $x>0$'' as its complete parameter
hypothesis, with no componentwise clause; in \cite{KS} the aggregate
condition is thus used as a substitute for $b_i>a_i$ and not as an addition
to it. Third, in the discussion following Theorem~3 the authors record that
\eqref{eq:bound} for $b_i\ge a_i>0$, $\sigma>0$ and $x>0$ is already due to
Luke, and that under the componentwise hypotheses the only case left open
is $0<\sigma<1$ with $-1<x<0$, which they conjecture there in prose; a
componentwise Conjecture~1 would therefore restate a theorem they credit to
Luke on $x>0$ and that prose conjecture on $-1<x<0$.

The two halves of Conjecture~1 have fared differently. The relaxation of
$\sigma\ge1$ to $\sigma>0$ is established on ranges controlled by
majorization, in \cite[Theorem~5]{KP} and, on a wider Meijer-$G$ range, in
\cite{Derbazi}, where the resulting lower bound is described as settling
``Conjecture~1 of \cite{KS} on that range'', with the remark that neither
bound obtained there ``addresses the broader conjectural condition
$\sum_i(b_i-a_i)>0$'' \cite[discussion following
Proposition~3.3]{Derbazi}. It is that broader condition, the replacement of
$b_i>a_i>0$ by \eqref{eq:aggregate}, which fails. The counterexample below
is taken at $x=1/2$, which lies in the range $x>-1$ of Theorem~3 as well as
in the range $x>0$ of the bound credited to Luke in \cite{KS}, so the
refutation does not depend on which of the two $x$-ranges is attached to
Conjecture~1.

\begin{theorem}\label{thm:counterexample}
The aggregate-condition form of Conjecture~1 of Karp and Sitnik is false.
In fact,
\[
 H_{2,(2,2),(1,5)}\!\left(\frac12\right)
 <
 \left(1+\frac45\cdot\frac12\right)^{-2}
 =\frac{25}{49}.
\]
\end{theorem}

\begin{proof}
The parameters satisfy
\[
 (1-2)+(5-2)=2>0,
 \qquad
 \mu((2,2),(1,5))=\frac45.
\]
The coefficient of $z^n$ in ${}_3F_2(2,2,2;1,5;z)$ is
\[
 \frac{(2)_n^3}{(1)_n(5)_n n!}
 =
 \frac{24(n+1)^2}{(n+2)(n+3)(n+4)}
 =
 \frac{12}{n+2}-\frac{96}{n+3}+\frac{108}{n+4}.
\]
For $0<|z|<1$ and an integer $r\ge1$,
\[
 \sum_{n=0}^{\infty}\frac{z^n}{n+r}
 =
 z^{-r}\left(
 -\log(1-z)-\sum_{j=1}^{r-1}\frac{z^j}{j}
 \right).
\]
Substituting $z=-1/2$ gives
\begin{equation}\label{eq:closed}
 {}_3F_2\!\left(
 \begin{matrix}2,2,2\\1,5\end{matrix};-\frac12
 \right)
 =
 1032-2544\log\frac32.
\end{equation}
Moreover,
\[
 \log\frac32
 =
 2\operatorname{artanh}\frac15
 =
 2\sum_{k=0}^{\infty}\frac{1}{(2k+1)5^{2k+1}}
 >
 2\left(\frac15+\frac1{375}+\frac1{15625}\right)
 =
 \frac{19006}{46875}.
\]
Therefore
\[
 1032-2544\log\frac32
 <
 \frac{7912}{15625}
 <
 \frac{25}{49},
\]
where the last inequality follows from
\[
 25\cdot15625-49\cdot7912=2937>0.
\]
\end{proof}

The value $q=2$ is minimal for this phenomenon, since for $q=1$ condition
\eqref{eq:aggregate} is exactly $b_1>a_1$.

\begin{proposition}\label{prop:local}
Let $A,B\in(0,\infty)^q$ and $\sigma>0$, and define
\[
 P_A:=\prod_{i=1}^q\left(1+\frac1{a_i}\right),
 \qquad
 P_B:=\prod_{i=1}^q\left(1+\frac1{b_i}\right).
\]
If $P_B>P_A$, then there exists $\varepsilon>0$ such that
\[
 H_{\sigma,A,B}(x)<(1+\mu(A,B)x)^{-\sigma}
 \qquad (0<x<\varepsilon).
\]
\end{proposition}

\begin{proof}
As $x\to0$,
\begin{align*}
 H_{\sigma,A,B}(x)
 &=
 1-\sigma\mu x
 +\frac{\sigma(\sigma+1)}2
  \mu^2\frac{P_A}{P_B}x^2
 +O(x^3),\\
 (1+\mu x)^{-\sigma}
 &=
 1-\sigma\mu x
 +\frac{\sigma(\sigma+1)}2\mu^2x^2
 +O(x^3),
\end{align*}
where $\mu=\mu(A,B)$. Hence
\[
 H_{\sigma,A,B}(x)-(1+\mu x)^{-\sigma}
 =
 \frac{\sigma(\sigma+1)}2\mu^2
 \left(\frac{P_A}{P_B}-1\right)x^2
 +O(x^3),
\]
whose leading coefficient is negative when $P_B>P_A$. The difference is
therefore negative for all sufficiently small $x>0$.
\end{proof}

\begin{corollary}\label{cor:family}
For every integer $m\ge2$ and every $\sigma>0$, let
\[
 A=(m,m),
 \qquad
 B=(1,m^2+1).
\]
Then \eqref{eq:aggregate} holds and $\mu(A,B)<1$, but \eqref{eq:bound} fails
for all sufficiently small $x>0$.
\end{corollary}

\begin{proof}
Here
\[
 \sum_{i=1}^2(b_i-a_i)=m^2-2m+2>0,
 \qquad
 \mu(A,B)=\frac{m^2}{m^2+1}<1.
\]
Also
\[
 P_A=\left(1+\frac1m\right)^2,
 \qquad
 P_B=2\left(1+\frac1{m^2+1}\right).
\]
For $m=2$,
\[
 P_A=\frac94<\frac{12}{5}=P_B;
\]
for $m\ge3$,
\[
 P_A\le\frac{16}{9}<2<P_B.
\]
Proposition~\ref{prop:local} applies.
\end{proof}

Karp and Prilepkina later proved \eqref{eq:bound} for every $\sigma>0$
under weak supermajorization \cite[Theorem~5]{KP}; there $B\prec^W A$ means
that, both vectors being arranged increasingly,
$\sum_{i=1}^{k}a_i\le\sum_{i=1}^{k}b_i$ for $k=1,\ldots,q$. The present
example is not covered: after increasing rearrangement the case $k=1$ of
this condition would require $a_1\le b_1$, that is, $2\le1$. Its case
$k=q$, on the other hand, is \eqref{eq:aggregate} with equality allowed, so
Theorem~\ref{thm:counterexample} shows that the hypothesis $B\prec^W A$ in
\cite[Theorem~5]{KP} cannot be weakened to that final partial sum alone.

\begin{thebibliography}{9}
\small

\bibitem{KS}
D. Karp and S. M. Sitnik,
\emph{Inequalities and monotonicity of ratios for generalized hypergeometric
function},
J. Approx. Theory \textbf{161} (2009), no.~1, 337--352.
\href{https://doi.org/10.1016/j.jat.2008.10.002}
{doi:10.1016/j.jat.2008.10.002};
\href{https://arxiv.org/abs/math/0703084}
{arXiv:math/0703084v2}.
The theorem and conjecture numbers cited above are those of the arXiv
version.

\bibitem{Derbazi}
Z. Derbazi,
\emph{A probabilistic sign rule for quotients of positive series and integral
transforms},
\href{https://arxiv.org/abs/2607.02511}
{arXiv:2607.02511v1 [math.CA]}, 2026.

\bibitem{KP}
D. Karp and E. Prilepkina,
\emph{Hypergeometric functions as generalized Stieltjes transforms},
J. Math. Anal. Appl. \textbf{393} (2012), no.~2, 348--359.
\href{https://doi.org/10.1016/j.jmaa.2012.03.044}
{doi:10.1016/j.jmaa.2012.03.044}.

\end{thebibliography}

\end{document}